\chapter{Brakke flows and loss of mass}\label{sec:brakke-flow}
A weak flow records time-dependent area measures and an inequality expressing dissipation. It may lose mass abruptly. We use an integrated, boundaryless formulation, stating the integrability and time-interval requirements explicitly. This makes the difference from the smooth identity in Proposition \ref{prop:mcf-test-transport} visible. Sources are \cite{WhiMCF15}, Section 9, and \cite{CCMS20}, Sections 2.4--2.6.

\begin{definition}[Integral Brakke flow without boundary]\label{def:integral-brakke-flow}
Let $U\subset\mathbb R^{n+1}$ be open and $I$ an interval. An integral Brakke flow in the formulation used here is a weakly measurable family of Radon measures $(\mu_t)_{t\in I}$ with the following properties. For almost every $t$, $\mu_t=\|V_t\|$ for an integral $n$-varifold, and its first variation has no singular part:
\[
 \delta V_t(X)=-\int\vec H_t\cdot X\,d\mu_t,
 \qquad \vec H_t\perp T_xV_t\quad\mu_t\text{-almost everywhere}.
\]
On each compact $K\subset U$ and compact interval $[a,b]\subset I$ we require uniform local mass bounds and
\[
 \sup_{t\in[a,b]}\mu_t(K)<\infty,
 \qquad \int_a^b\int_K|\vec H_t|^2\,d\mu_t\,dt<\infty.
\]
For every such $[a,b]$ and every nonnegative $\phi\in C_c^1(U\times[a,b])$,
\[
 \mu_b(\phi(\cdot,b))-\mu_a(\phi(\cdot,a))
 \leq\int_a^b\int\left(\partial_t\phi+\nabla\phi\cdot\vec H_t
 -\phi|\vec H_t|^2\right)d\mu_t\,dt.
\]
Here $\mu_t(f)=\int f\,d\mu_t$. The integral inequality is imposed on every subinterval, not only at almost every pair of endpoints. The mass and curvature bounds are listed as part of the working formulation; compare \cite{WhiMCF15}, Definition 9.1 and Corollary 9.4--Theorem 9.5. Prescribed boundary terms are not included in this definition.
\end{definition}

\begin{proposition}[A smooth closed flow is a Brakke flow]\label{prop:smooth-to-brakke}
The area measures of a smooth closed embedded mean curvature flow satisfy Definition \ref{def:integral-brakke-flow}. Its integrated inequality holds with equality.
\end{proposition}
\begin{proof}
Each smooth time slice gives a multiplicity-one integral varifold, whose first variation is $-\int\vec H\cdot X$. Smoothness on the compact parameter manifold and a compact time interval gives both local mass and curvature bounds. Integrate Proposition \ref{prop:mcf-test-transport} from $a$ to $b$. The same argument is local for test functions supported away from a boundary.
\end{proof}

\begin{example}[A plane that suddenly disappears]\label{ex:brakke-disappearing-plane}
Let $P$ be an affine $n$-plane, let $q\geq1$ be an integer, and fix a time $t_*$. Set $\mu_t=q\mathcal H^n\llcorner P$ for $t<t_*$ and $\mu_t=0$ for $t\geq t_*$. This is an integral Brakke flow. To verify the inequality, note that $\vec H=0$ whenever the measure is nonzero. On an interval crossing $t_*$, the time integral of $\partial_t\phi$ is
\[
 q\int_P\phi(x,t_*)\,d\mathcal H^n-q\int_P\phi(x,a)\,d\mathcal H^n,
\]
while the difference of the endpoint measures is just the negative second term. Nonnegativity of $\phi$ gives the required inequality. On intervals not crossing $t_*$ it is an equality. Thus the Brakke inequality alone permits disappearance even though curvature was zero. Compare \cite{WhiMCF15}, Remark 9.2.
\end{example}

\begin{proposition}[A shrinking spatial cutoff]\label{prop:brakke-shrinking-cutoff}
Suppose $\overline{B_R(a)}\subset U$ and $s\leq t$ are flow times. Put
\[
 q(x,t)=\bigl(R^2-|x-a|^2-2n(t-s)\bigr)_+,
 \qquad \phi(x,t)=q(x,t)^3.
\]
Then $t\mapsto\mu_t(\phi(\cdot,t))$ is nonincreasing for $t\geq s$. In particular, for $0<r<R$ and $0\leq t-s\leq(R^2-r^2)/(4n)$,
\[
 \mu_t(B_r(a))\leq
 \left(\frac{2R^2}{R^2-r^2}\right)^3\mu_s(B_R(a)).
\]
This is a cubic-cutoff version of the local mass estimate in \cite{WhiMCF15}, Theorem 9.3.
\end{proposition}
\begin{proof}
The cube of the positive part is $C^2$, with compact spatial support in $U$ for $t\geq s$. On the region where $q>0$, and for any $n$-plane $S$,
\[
 \partial_t\phi=-6nq^2,\qquad
 \operatorname{tr}_S D^2\phi=6q|\nabla^Sq|^2-6nq^2.
\]
Thus $\partial_t\phi-\operatorname{tr}_SD^2\phi\leq0$, also across $q=0$. The first-variation identity with $X=\nabla\phi$ gives $\int\nabla\phi\cdot\vec H_t\,d\mu_t=-\int\operatorname{tr}_{T_xV_t}D^2\phi\,d\mu_t$ at almost every time. Insert this in the integrated Brakke inequality and discard $-\phi|\vec H|^2$. This proves monotonicity. On $B_r(a)$ at the stated times, $q\geq(R^2-r^2)/2$, whereas initially $\phi\leq R^6$. These two bounds give the displayed mass estimate.
\end{proof}

\begin{proposition}[Local curvature energy from the Brakke inequality]\label{prop:brakke-local-energy}
Let $\eta\in C_c^1(U)$ and $[a,b]\subset I$. Then
\[
 \int_a^b\int\eta^2|\vec H_t|^2\,d\mu_t\,dt
 \leq2\mu_a(\eta^2)-2\mu_b(\eta^2)
 +4\int_a^b\int|D\eta|^2\,d\mu_t\,dt.
\]
The last integral is controlled by local mass bounds on the support of $\eta$. This is the energy estimate of \cite{WhiMCF15}, Theorem 9.5, written with a squared cutoff.
\end{proposition}
\begin{proof}
Use the time-independent test function $\phi=\eta^2$. Young's inequality gives
$2\eta\nabla\eta\cdot\vec H\leq\frac12\eta^2|\vec H|^2+2|\nabla\eta|^2$.
Move the remaining negative curvature term to the other side of the Brakke inequality and multiply by two. In particular a cutoff equal to one on a compact set bounds the spacetime curvature energy there.
\end{proof}

\begin{proposition}[One-sided limits of localized mass]\label{prop:brakke-time-semicontinuity}
For $\eta\in C_c^1(U)$ and a compact time interval $[a,b]\subset I$, there is $C<\infty$ such that
\[
 t\longmapsto\mu_t(\eta^2)-Ct
\]
is nonincreasing on $[a,b]$. Consequently its left and right limits exist at interior times, and
\[
 \lim_{s\uparrow t}\mu_s(\eta^2)\geq\mu_t(\eta^2)
 \geq\lim_{s\downarrow t}\mu_s(\eta^2).
\]
These inequalities allow downward jumps. They are not continuity in time. Compare \cite{WhiMCF15}, Theorems 9.6--9.7.
\end{proposition}
\begin{proof}
Complete the square in the time-independent Brakke integrand:
\[
 -\eta^2|\vec H|^2+2\eta\nabla\eta\cdot\vec H
 =-|\eta\vec H-\nabla\eta|^2+|\nabla\eta|^2\leq|\nabla\eta|^2.
\]
Let $M$ bound $\mu_t(\operatorname{spt}\eta)$ on the interval, and take $C=M\|D\eta\|_\infty^2$. The integrated inequality between arbitrary times gives the monotonicity. The one-sided conclusions follow from elementary properties of a monotone real function.
\end{proof}

\begin{definition}[Convergence of integral Brakke flows]\label{def:brakke-convergence}
We say $\mathcal M_j=(\mu^j_t)$ converges to $\mathcal M=(\mu_t)$ on a common space-time domain if $\mu^j_t\rightharpoonup\mu_t$ as Radon measures at every time, and for almost every time one can pass to a further subsequence, possibly depending on the time, for which the associated integral varifolds converge to the varifold of the limiting flow. This is the convention of \cite{CCMS20}, Section 2.6. The time-dependent subsequence in the second clause is not a single subsequence with varifold convergence at every time.
\end{definition}

\begin{theorem}[Compactness of integral Brakke flows]\label{thm:brakke-compactness}
Let $\mathcal M_j$ be integral Brakke flows on $U\times[a,b]$. Suppose for every compact $K\subset U$,
\[
 \sup_j\sup_{t\in[a,b]}\mu_t^j(K)<\infty.
\]
Then a subsequence converges to an integral Brakke flow in the sense of Definition \ref{def:brakke-convergence}. Source: \cite{WhiMCF15}, Theorem 9.8, pp.~23--24; compare \cite{CCMS20}, Section 2.6. The proof combines the localized mass monotonicity, curvature energy bounds, varifold compactness, and lower semicontinuity of the Brakke functional. That proof remains external. This is weak compactness and does not assert smooth convergence.
\end{theorem}

\begin{theorem}[Gaussian monotonicity for Brakke flows]\label{thm:brakke-gaussian-monotonicity}
Let $\mathcal M$ be an integral Brakke flow on $\mathbb R^{n+1}\times I$. Assume that on each compact time interval its measures have a uniform entropy bound. For $[a,b]\subset I$ and any $t_0>b$,
\[
 \int\rho_{x_0,t_0}(x,b)\,d\mu_b-
 \int\rho_{x_0,t_0}(x,a)\,d\mu_a
 \leq-\int_a^b\int\left|\vec H_t+
 \frac{(x-x_0)^\perp}{2(t_0-t)}\right|^2\rho_{x_0,t_0}\,d\mu_t\,dt.
\]
The Gaussian ratios and density are then defined as in Corollary \ref{cor:mcf-gaussian-density-limit}. Entropy is nonincreasing. The entropy bound supplies spatial tail control for using a Gaussian in place of compactly supported test functions; compare Proposition \ref{prop:mcf-entropy-area-growth}. Sources: \cite{CCMS20}, Section 2.5, and \cite{WhiMCF15}, Section 10.1. The cutoff and weak lower-semicontinuity argument remains external; equality is not claimed.
\end{theorem}

\begin{definition}[Unit-regular Brakke flows]\label{def:unit-regular-brakke}
An integral Brakke flow is unit-regular if every interior spacetime point with Gaussian density one has an open spacetime neighborhood on which the flow is a smooth embedded multiplicity-one mean curvature flow. This is a two-sided neighborhood in time. It is an additional property, not part of Definition \ref{def:integral-brakke-flow}. See \cite{CCMS20}, Section 2.6.
\end{definition}

\begin{example}[Density one does not prevent a mass drop]\label{ex:brakke-density-one-drop}
Take the disappearing plane of Example \ref{ex:brakke-disappearing-plane} with $q=1$. At $(x,t_*)$ with $x\in P$, every backward Gaussian ratio is one, so the Gaussian density is one. Nevertheless every two-sided spacetime neighborhood meets the vanished part of the flow. Thus this flow is not unit-regular. A backward regularity statement and a two-sided regularity statement are different at a mass-drop time; compare \cite{WhiMCF15}, the discussion after Theorem 11.2.
\end{example}

\begin{definition}[Cyclic integral Brakke flows]\label{def:cyclic-brakke}
An integral Brakke flow is cyclic modulo two if at almost every time its associated rectifiable mod-two flat chain satisfies
\[
 \partial[V_t]=0.
\]
Here $[V_t]$ reduces the integer multiplicity modulo two, as in Definition \ref{def:varifold-mod-two}. This property concerns a flat-chain boundary and does not assert that the varifold is orientable or that the flow is unit-regular. See \cite{CCMS20}, Section 2.6, and \cite{White09}.
\end{definition}
